% abstract.tex -- Sequencer Paper (BNR-2026-009)

\begin{abstract}
Layer~2 sequencer design presents a fundamental tension between operator efficiency and
censorship resistance: a single-operator sequencer achieves high throughput and low latency
but introduces a trust assumption that the operator will include all valid transactions
fairly. We investigate the design and performance of a single-operator sequencer for Basis
Network, an enterprise zkEVM L2 deployed on Avalanche where each enterprise operates its
own chain. Through systematic analysis of eight production L2 systems---zkSync Era,
Polygon CDK, Scroll, Arbitrum One, Optimism (OP Stack), Taiko, MegaETH, and
Starknet---we identify universal patterns: all production rollups use centralized sequencers,
and forced inclusion mechanisms via L1 provide the censorship resistance guarantee.
We implement a Go prototype (${\sim}$500~LOC) featuring a FIFO mempool, Arbitrum-style
forced inclusion queue with configurable deadlines, and a tick-based block production
loop. Benchmarks across 30+ repetitions per configuration demonstrate block production
latency of $0.019$~ms at 100~transactions per block scaling to $0.885$~ms at
5{,}000~transactions, mempool insertion throughput of ${\sim}$2.8~million~tx/s
(365~ns/op), and 100\% FIFO ordering accuracy across all test scenarios including
concurrent access with 4~producer goroutines at 33{,}000~tx/s aggregate insertion rate.
Forced inclusion achieves 100\% transaction inclusion within one block tick
(${\sim}$1~second) under both cooperative and adversarial conditions. These results
confirm that a single-operator sequencer can sustain 1-second block times under enterprise
workloads while an L1-based forced inclusion mechanism guarantees censorship resistance
with bounded latency. Five architectural invariants (INV-SEQ-1 through INV-SEQ-5) are
established as constraints for downstream formalization, implementation, and verification.
\end{abstract}
